\chapter{The theorems, side by side: human proof and Lean~4}\label{ch2:chapter}

This chapter states every result of the book twice. First comes the mathematics: an informal statement, a precise
statement and a human proof that you can check with pen and paper. Then comes the Lean~4 text, quoted verbatim from the
files in \lean{formal_cert/Triangular/}, together with a short paragraph explaining how the formal proof follows the
human one. Where a listing is shortened, the omission is marked by a comment \lean{-- ...}; nothing else in a listing is
changed.

\paragraph{How the tiers are used in this chapter.}
\begin{itemize}
\item A Lean theorem that compiled under the pinned toolchain (Lean~4.34.1, Mathlib \lean{d13f23b7}, openai/math
  \lean{adc7f124}), with an axiom footprint inside \{\lean{propext}, \lean{Classical.choice}, \lean{Quot.sound}\}, is
  \tierA. Tier~A attaches to the \emph{formal statement}. It means ``kernel-checked, conditional on the fidelity of the
  definitions''. It does not mean ``true of the physical world'', and it does not mean that the Lean definitions say what
  the English says: no human has audited the adequacy of the statements, and the claims ledger carries six open
  adequacy flags.
\item Every result below except the upstream theorem is a corollary of the upstream theorem
  \lean{OAI.AtomicTriangular.universal_energy_minimum}~\cite{oai}. That theorem is AI-generated (OpenAI release). We
  kernel-checked it and Comparator accepted it under the pinned toolchain. It has not been peer reviewed.
\item A human proof written out in this book is an argument, \tierC, until it is matched by a kernel-checked Lean
  proof. When it is matched, the \emph{Lean statement} is Tier~A; the human proof stays an explanation of it.
\item Statements taken from the upstream manuscript or the published literature are \tierL.
\end{itemize}

\begin{table}[ht]
\centering\footnotesize
\begin{tabular}{@{}l p{9.4cm} p{3.6cm}@{}}
\toprule
Result & Main Lean name (file) & Tier \\
\midrule
Upstream & \lean{universal_energy_minimum} (upstream) & A (AI-generated, not peer reviewed) \\
D1 & \lean{triangular_riesz_optimal} (\lean{Riesz.lean}) & A \\
D1b & \lean{latticeEnergy_riesz_lt_top} (\lean{RieszFinite.lean}) & A \\
D2 & \lean{triangular_optimal_any_density} (\lean{AnyDensity.lean}) & A \\
Controls & \lean{riesz_neg_two_not_admissible}, \dots{} (\lean{Controls.lean}) & A \\
D3 & \lean{universal_optimum_not_unique} (\lean{UniversalOptimality.lean}) & A \\
D4 & \lean{energy_eq_openEnergy}, \lean{densityOne_iff_open} (\lean{OpenDisk.lean}, \lean{OpenDiskAny.lean}) & A \\
\bottomrule
\end{tabular}
\caption{The results of this chapter. ``A'' always means: kernel-checked, conditional on the fidelity of the
definitions.}\label{ch2:tab:results}
\end{table}

% =====================================================================================================
\section{The objects, as upstream defines them}\label{ch2:sec:defs}

We write $\mathbb{P}=\R^2$ with the Euclidean norm, and $B_R=\{x:\|x\|\le R\}$ for the \emph{closed} disk of radius
$R$ centred at the origin. All definitions in this section are upstream's~\cite{oai}. They are taken from the challenge
file \lean{lean/ComparatorChallenges/TriangularEnergy.lean} (openai/math \lean{adc7f124}, sha256 \lean{af51979c...0616};
we checked the full hash, \lean{af51979c7baecd3e9852637bbba29eff}\allowbreak\lean{96cbb09cec84e15eff7d686f1a770616}, with
\lean{sha256sum} on our pinned copy).

\begin{definition}[Configurations, counts, density]\label{ch2:def:basic}
A \emph{configuration} is a set $C\subseteq\mathbb{P}$ (no multiplicities). It is \emph{locally finite} if
$C\cap B_R$ is finite for every $R$. Its \emph{disk count} is $N_R(C)=\#(C\cap B_R)$; by convention the count is $0$
if $C\cap B_R$ is infinite. $C$ has \emph{density one} if $N_R(C)/(\pi R^2)\to 1$ as $R\to\infty$, and density $\rho$
if the limit is $\rho$.
\end{definition}

\begin{definition}[Admissible potentials]\label{ch2:def:adm}
$g:\R\to\R$ is \emph{admissible} if $g$ is $C^\infty$ on $(0,\infty)$, $g(t)\ge 0$ for $t>0$, and
$(-1)^r g^{(r)}(t)\ge 0$ for every $r\in\N$ and every $t>0$. The inequalities are \emph{weak}, and the smoothness is
$C^\infty$, not analyticity. The variable $t$ is the \emph{squared} distance: a pair at distance $r$ interacts through
$g(r^2)$.
\end{definition}

\begin{definition}[Disk energy, lower energy, lattice energy]\label{ch2:def:energy}
With values in $[0,\infty]$,
\[
E_g(C;R)=\frac{1}{N_R(C)}\sum_{x\in C\cap B_R}\ \sum_{\substack{y\in C\cap B_R\\ y\neq x}} g\bigl(\|x-y\|^2\bigr),
\qquad
E_g(C)=\liminf_{R\to\infty}E_g(C;R).
\]
(Each term is truncated at $0$ by \lean{ENNReal.ofReal}, and $0^{-1}=\infty$, $\infty\cdot 0=0$ in $[0,\infty]$, so an
empty disk has energy $0$.) $E_g(C)$ is the \emph{lower} energy: it is a $\liminf$, and nothing requires the limit to
exist. Let $b=\sqrt3/2$ and
\[
p(j,k)=b^{-1/2}\Bigl(j+\tfrac k2,\ kb\Bigr),\qquad A=\{p(j,k):(j,k)\in\Z^2\}.
\]
$A$ is the triangular lattice of covolume one. Its \emph{lattice energy} is $\mathcal{E}_g(A)=\sum_{a\in A\setminus\{0\}}
g(\|a\|^2)\in[0,\infty]$; the sum may be $\infty$.
\end{definition}

\begin{leanbox}[(upstream challenge file, verbatim)]
\begin{lstlisting}[language=Lean,escapeinside={(*}{*)}]
-- lean/ComparatorChallenges/TriangularEnergy.lean, lines 11-36 (namespace OAI.TriangularUniversal)
abbrev Plane := EuclideanSpace ℝ (Fin 2)

def LocallyFinite (C : Set Plane) : Prop :=
  ∀ R : ℝ, (C ∩ Metric.closedBall (0 : Plane) R).Finite

def diskPoints (C : Set Plane) (R : ℝ) : Finset Plane :=
  if h : (C ∩ Metric.closedBall (0 : Plane) R).Finite then h.toFinset else (*$\emptyset$*)

def diskCount (C : Set Plane) (R : ℝ) : ℕ := (diskPoints C R).card

def DensityOne (C : Set Plane) : Prop :=
  Tendsto (fun R : ℝ => (diskCount C R : ℝ) / (Real.pi * R ^ 2))
    atTop (𝓝 1)

def AdmissiblePotential (g : ℝ → ℝ) : Prop :=
  ContDiffOn ℝ (⊤ : ℕ∞) g (Set.Ioi 0) ∧
  (∀ t : ℝ, 0 < t → 0 ≤ g t) ∧
  ∀ (r : ℕ) (t : ℝ), 0 < t → 0 ≤ (-1 : ℝ) ^ r * iteratedDeriv r g t

def diskEnergy (g : ℝ → ℝ) (C : Set Plane) (R : ℝ) : ℝ≥0∞ :=
  (diskCount C R : ℝ≥0∞)⁻¹ *
    ∑ x ∈ diskPoints C R, ∑ y ∈ (diskPoints C R).erase x,
      ENNReal.ofReal (g (‖x - y‖ ^ 2))

def energy (g : ℝ → ℝ) (C : Set Plane) : ℝ≥0∞ :=
  Filter.liminf (diskEnergy g C) atTop
-- ... lines 40-63 (namespace OAI.AtomicTriangular)
def b : ℝ := Real.sqrt 3 / 2

def triangularPoint (jk : ℤ × ℤ) : Plane :=
  (Real.sqrt b)⁻¹ •
    (EuclideanSpace.single 0 ((jk.1 : ℝ) + (jk.2 : ℝ) / 2) +
      EuclideanSpace.single 1 ((jk.2 : ℝ) * b))

def A : Set Plane := Set.range triangularPoint
-- ...
def latticeEnergy (g : ℝ → ℝ) : ℝ≥0∞ :=
  ∑' a : {x : Plane // x ∈ A ∧ x ≠ 0}, ENNReal.ofReal (g (‖a.val‖ ^ 2))
\end{lstlisting}
\end{leanbox}

Two details of the formal definitions are worth noticing before any proof. First, \lean{diskPoints} returns the empty
set when $C\cap B_R$ is infinite, so every useful statement carries the hypothesis \lean{LocallyFinite C}. Second,
\lean{iteratedDeriv} is the iterated \emph{global} derivative of $g:\R\to\R$; at a point $t>0$ it agrees with the
classical derivative because $g$ is smooth on the open set $(0,\infty)$ around $t$ (this is item ``Potential class'' of
the statement-fidelity audit, \lean{results/discovery/statement_fidelity.md}) \tierC.

% =====================================================================================================
\section{Result 0: the upstream theorem}\label{ch2:sec:upstream}

\paragraph{Informal statement.} Among all locally finite planar configurations of density one, the triangular lattice
has the least lower energy, simultaneously for every admissible potential.

\begin{theorem}[Universal energy minimum~\cite{oai}]\label{ch2:thm:upstream}
Let $g$ be admissible and $C$ locally finite of density one. Then
\[
\mathcal{E}_g(A)\le E_g(C)\qquad\text{and}\qquad \mathcal{E}_g(A)=E_g(A).
\]
\end{theorem}
Status: \tierA{} for the Lean statement below (AI-generated proof, kernel-checked and Comparator-accepted by us under
the pinned toolchain, footprint \{\lean{propext}, \lean{Classical.choice}, \lean{Quot.sound}\}; not peer reviewed);
\tierL{} for the manuscript's proof.

\paragraph{What is claimed, and what is not.}
\begin{itemize}
\item The theorem is an \emph{energy comparison}: it identifies the minimum \emph{value} $\mathcal{E}_g(A)$ of the lower
  energy. In upstream's words, ``It identifies the minimum value, not all minimizers.'' \tierL{} Section~\ref{ch2:sec:d3}
  shows, in Lean, that the minimizer is indeed not unique.
\item Upstream does \emph{not} claim the per-potential form of the Cohn--Kumar conjecture, in which each potential has a
  sharp auxiliary function of its own. That strong form is \textbf{open}. \tierL
\item The competitor $C$ need not be periodic; it only has to be locally finite with density one.
\item Energies are lower energies over \emph{closed} disks centred at $0$. Section~\ref{ch2:sec:d4} shows that open disks
  give the same density and the same lower energy.
\item The Riesz endpoint $s=2$ and renormalized (Coulomb-type) energies are not covered. Renormalized energies are
  treated only in the upstream manuscript (its Section~8), and not in Lean anywhere. \tierL
\end{itemize}

\subsection{Proof architecture (Tier L)}
We give only the architecture of upstream's proof (upstream manuscript, Section~1), not its details. Everything in this
subsection is \tierL.

\paragraph{Step 1: Gaussians.} Use the normalisation $\hat f(\xi)=\int f(x)e^{-2\pi i x\cdot\xi}\,dx$, the Gaussian
$G_\alpha(x)=e^{-\pi\alpha\|x\|^2}$ ($\alpha>0$) and the dual lattice $A^*$. For every $\alpha>0$ upstream constructs a
radial Schwartz function $f_\alpha$ with
\[
f_\alpha\le G_\alpha,\qquad \hat f_\alpha\ge 0,\qquad f_\alpha(a)=G_\alpha(a)\ (a\in A\setminus\{0\}),\qquad
\hat f_\alpha(w)=0\ (w\in A^*\setminus\{0\}).
\]
So $f_\alpha$ is a \emph{minorant} of the potential (not a majorant), with nonnegative Fourier transform.

\paragraph{Step 2: the linear-programming bound.} For every competitor $C$ of centred density one,
\[
E_{G_\alpha}(C)\ \ge\ \hat f_\alpha(0)-f_\alpha(0).
\]
This is a Cohn--Elkies-type bound~\cite{ck,ckmrv}. The competitor does not need to be periodic; upstream compares the
disk sums with the area of the disk and controls the error uniformly.

\paragraph{Step 3: equality on the lattice.} Poisson summation on $A$ (covolume one) reads
$\sum_{a\in A}f_\alpha(a)=\sum_{w\in A^*}\hat f_\alpha(w)$. The interpolation conditions turn it into
$f_\alpha(0)+\sum_{a\neq0}G_\alpha(a)=\hat f_\alpha(0)$. Hence
\[
\hat f_\alpha(0)-f_\alpha(0)=\sum_{a\in A\setminus\{0\}}e^{-\pi\alpha\|a\|^2}=\mathcal{E}_{G_\alpha}(A),
\]
and the bound of Step~2 is attained by $A$.

\paragraph{Step 4: the construction of $f_\alpha$.} A sine product with double zeros encodes the periodic contact set.
Multiplying it by double-pole and simple-pole terms gives entire functions with prescribed values and slopes.
Compactly supported spectral measures and Gaussian damping follow. Finite reference lists are handled by positive
quadrature, with certified matrix and polynomial inequalities (the verification scripts use \lean{python-flint} ball
arithmetic). The construction does \emph{not} use modular forms; modular forms are the method of~\cite{ckmrv} in
dimensions $8$ and $24$.

\paragraph{Step 5: from Gaussians to all admissible $g$.} An admissible $g$ is a positive mixture of Gaussians in
$t=\|x\|^2$: $g(t)=\int_{[0,\infty)}e^{-\pi\alpha t}\,d\mu(\alpha)$, i.e.\ $g(\|x\|^2)=\int G_\alpha(x)\,d\mu(\alpha)$
in the normalisation of Step~1 (Bernstein--Hausdorff--Widder, after rescaling $\alpha$ by $\pi$; upstream gives its own
finite-difference proof). The measure may charge $\alpha=0$, i.e.\ contain a constant, whose energy is infinite at
density one (upstream lemma \lean{energy_one}). Integrating the Gaussian bounds needs monotone convergence and Tonelli,
and, because $\liminf$ is not linear, a separate lower-limit argument; the endpoint and infinite-energy cases are treated
in upstream's Section~5.

\begin{leanbox}[(upstream proof, top level, verbatim)]
The challenge file states the theorem with \lean{sorry} by design; the proof lives in
\lean{OAI/Analysis/Triangular/Energy/Universal.lean}, and Comparator certifies that the two statements agree. We quote
the 26-line proof file from our copy of the upstream release \lean{fd4aeeb2} (the triangular files are unchanged
from \lean{adc7f124}).
\begin{lstlisting}[language=Lean,escapeinside={(*}{*)}]
-- ComparatorChallenges/TriangularEnergy.lean, lines 65-69
/-- Universal energy minimality for smooth completely monotone potentials, with attainment. -/
theorem universal_energy_minimum (g : ℝ → ℝ) (C : Set Plane)
    (hg : AdmissiblePotential g) (hC : LocallyFinite C) (hd : DensityOne C) :
    latticeEnergy g ≤ energy g C ∧ latticeEnergy g = energy g A := by
  sorry

-- OAI/Analysis/Triangular/Energy/Universal.lean, lines 7-22
namespace AtomicTriangular
/-- The triangular lattice minimizes every nonnegative completely monotone energy. -/
theorem universal_energy_minimum (g : ℝ → ℝ) (C : Set Plane)
    (hg : AdmissiblePotential g) (_hC : LocallyFinite C) (hd : DensityOne C) :
    latticeEnergy g ≤ energy g C ∧ latticeEnergy g = energy g A := by
  have hb (α : ℝ) (hα : 0<α) :
      TriangularUniversal.latticeEnergy (TriangularUniversal.gaussPot α)≤
        TriangularUniversal.energy (TriangularUniversal.gaussPot α) C := by
    rw [(*$\leftarrow$*)latticeEnergy_eq_support]
    exact gaussian_energy_bound hd α hα
  have hu := TriangularUniversal.universal_from_gaussian_bounds hd hb hg
  constructor
  · rw [latticeEnergy_eq_support]
    exact hu.1
  · rw [latticeEnergy_eq_support,A_eq_support]
    exact hu.2
\end{lstlisting}
\end{leanbox}

\paragraph{How the Lean proof mirrors the architecture.} The top-level proof has exactly two moves. The local fact
\lean{hb} is Steps~1--4: a lower bound for every Gaussian potential \lean{gaussPot α}, obtained from
\lean{gaussian_energy_bound}. The call \lean{universal_from_gaussian_bounds} is Step~5: it passes from Gaussian bounds
to every admissible $g$ and also yields attainment. The rewrites \lean{latticeEnergy_eq_support} and \lean{A_eq_support}
translate between the challenge file's description of the lattice and the one used inside the library. Notice that the
hypothesis \lean{_hC : LocallyFinite C} is unused in the proof (the leading underscore says so): local finiteness is
part of the \emph{statement} that Comparator checks, and it matters for the meaning of \lean{diskPoints}, but the bound
itself does not need it \tierC{} (the hypothesis is unused in upstream's proof term; no theorem without it is stated).

% =====================================================================================================
\section{Result D1: Riesz potentials}\label{ch2:sec:d1}

\paragraph{Informal statement.} For every $s>0$ the triangular lattice minimizes the Riesz energy
$\sum|x-y|^{-s}$ among locally finite configurations of density one. The statement has content only for $s>2$.

The Riesz kernel at distance $r$ is $r^{-s}$, so in the variable $t=r^2$ the potential is
\[
g_s(t)=t^{-s/2}.
\]
In Lean this is \lean{TriangularRiesz.riesz s := fun t => t ^ (-(s / 2))}, and the control \lean{riesz_sq} checks
$g_s(r^2)=r^{-s}$ for $r>0$ \tierA.

\begin{lemma}[Admissibility]\label{ch2:lem:rieszadm}
For every $s>0$, $g_s$ is admissible.
\end{lemma}

\begin{humanproof}
Write $\sigma=s/2>0$. We check the three conditions of Definition~\ref{ch2:def:adm}.

\emph{Smoothness.} On $(0,\infty)$, $t\mapsto t^{-\sigma}=e^{-\sigma\log t}$ is a composition of smooth functions.

\emph{Positivity.} For $t>0$, $t^{-\sigma}>0$.

\emph{Complete monotonicity.} Differentiating $r$ times,
\[
g_s^{(r)}(t)=(-\sigma)(-\sigma-1)\cdots(-\sigma-r+1)\,t^{-\sigma-r}.
\]
The coefficient is the \emph{descending} Pochhammer symbol $(x)_r^{\downarrow}=x(x-1)\cdots(x-r+1)$ at $x=-\sigma$.
Pulling a factor $-1$ out of each of the $r$ factors gives
\[
(-\sigma)_r^{\downarrow}=(-1)^r\,\sigma(\sigma+1)\cdots(\sigma+r-1)=(-1)^r(\sigma)_r^{\uparrow},
\]
with the \emph{ascending} Pochhammer symbol $(\sigma)_r^{\uparrow}$. Therefore
\[
(-1)^r g_s^{(r)}(t)=(-1)^r(-1)^r(\sigma)_r^{\uparrow}\,t^{-\sigma-r}=(\sigma)_r^{\uparrow}\,t^{-\sigma-r}.
\]
Each factor $\sigma+i$ ($0\le i<r$) is positive because $\sigma>0$, so $(\sigma)_r^{\uparrow}>0$, and $t^{-\sigma-r}>0$.
The product is positive, in particular $\ge0$.
\end{humanproof}

\begin{theorem}[D1, Riesz optimality]\label{ch2:thm:d1}
For every $s>0$ and every locally finite $C$ of density one, $\sum_{a\in A\setminus\{0\}}\|a\|^{-s}\le E_{g_s}(C)$ in
$[0,\infty]$.
\end{theorem}

\begin{humanproof}
Apply Theorem~\ref{ch2:thm:upstream} with $g=g_s$, which is admissible by Lemma~\ref{ch2:lem:rieszadm}, and keep the first
conjunct.
\end{humanproof}

Status: \tierA{} for the Lean statement (conditional on upstream and on the fidelity of the definitions). The upstream
manuscript already says informally that the theorem ``includes every inverse power $g(t)=t^{-p}$ with $p>0$''; D1 is
therefore the kernel-checked form of a stated special case, not a new mathematical claim.

\begin{warnbox}[(informative only for $s>2$)]
For $0<s\le2$ the left-hand side is $\infty$, so Theorem~\ref{ch2:thm:d1} only says $E_{g_s}(C)=\infty$.
The reason is a counting argument \tierC: the lattice has about $2\pi k$ points with $k\le\|a\|<k+1$, each contributing
about $k^{-s}$, and $\sum_k k^{1-s}=\infty$ for $s\le2$ (Exercise~\ref{ch2:ex:diverge}). We have \emph{not} proved this
divergence in Lean. For $s>2$ the left-hand side is finite (D1b, \S\ref{ch2:sec:d1b}), so the theorem is not the
trivial $\infty\le\infty$. The endpoint $s=2$ is open.
\end{warnbox}

\begin{leanbox}[(\lean{formal_cert/Triangular/Riesz.lean}, verbatim)]
\begin{lstlisting}[language=Lean,escapeinside={(*}{*)}]
/-- Riesz potential in the variable `t = r²`. -/
noncomputable def riesz (s : ℝ) : ℝ → ℝ := fun t => t ^ (-(s / 2))

theorem riesz_admissible {s : ℝ} (hs : 0 < s) : AdmissiblePotential (riesz s) := by
  refine ⟨?_, ?_, ?_⟩
  · intro x hx
    exact (Real.contDiffAt_rpow_const_of_ne (ne_of_gt hx)).contDiffWithinAt
  · intro t ht
    exact (Real.rpow_pos_of_pos ht _).le
  · intro r t ht
    have ha : 0 < s / 2 := by positivity
    rw [iteratedDeriv_eq_iterate]
    unfold riesz
    rw [Real.iter_deriv_rpow_const]
    have hasc := ascPochhammer_eval_neg_eq_descPochhammer (R := ℝ) (-(s / 2)) r
    rw [neg_neg] at hasc
    have hdesc : (descPochhammer ℝ r).eval (-(s / 2)) = (-1) ^ r * (ascPochhammer ℝ r).eval (s / 2) := by
      have h1 : ((-1 : ℝ) ^ r) * ((-1) ^ r) = 1 := by rw [(*$\leftarrow$*) mul_pow]; norm_num
      calc (descPochhammer ℝ r).eval (-(s / 2))
          = ((-1) ^ r * (-1) ^ r) * (descPochhammer ℝ r).eval (-(s / 2)) := by rw [h1, one_mul]
        _ = (-1) ^ r * (ascPochhammer ℝ r).eval (s / 2) := by rw [hasc]; ring
    rw [hdesc]
    have hpos := ascPochhammer_pos r (s / 2) ha
    have hpow := Real.rpow_pos_of_pos ht (-(s / 2) - r)
    have h1 : ((-1 : ℝ) ^ r) * ((-1) ^ r) = 1 := by rw [(*$\leftarrow$*) mul_pow]; norm_num
    calc (0 : ℝ) ≤ (ascPochhammer ℝ r).eval (s / 2) * t ^ (-(s / 2) - r) := by positivity
      _ = (-1) ^ r * ((-1) ^ r * (ascPochhammer ℝ r).eval (s / 2) * t ^ (-(s / 2) - r)) := by
        rw [(*$\leftarrow$*) mul_assoc, (*$\leftarrow$*) mul_assoc, h1, one_mul]

theorem triangular_riesz_optimal {s : ℝ} (hs : 0 < s) (C : Set Plane) (hC : LocallyFinite C) (hd : DensityOne C) :
    latticeEnergy (riesz s) ≤ energy (riesz s) C :=
  (universal_energy_minimum (riesz s) C (riesz_admissible hs) hC hd).1
\end{lstlisting}
\end{leanbox}

\paragraph{How the Lean proof mirrors the human proof.} The anonymous constructor \lean{⟨?_, ?_, ?_⟩} opens the three
conditions in the order of the human proof. Smoothness is the Mathlib fact that $x\mapsto x^p$ is smooth away from
$0$ (\lean{Real.contDiffAt_rpow_const_of_ne}); positivity is \lean{Real.rpow_pos_of_pos}. For complete monotonicity,
\lean{iteratedDeriv_eq_iterate} turns the $r$-th derivative into the $r$-fold iterate of \lean{deriv}, and
\lean{Real.iter_deriv_rpow_const} is the displayed formula with the descending Pochhammer polynomial evaluated at
$-\sigma$. The identity $(-\sigma)^{\downarrow}_r=(-1)^r(\sigma)^{\uparrow}_r$ is \lean{ascPochhammer_eval_neg_eq_descPochhammer}
read backwards; it is packaged as \lean{hdesc}, using $(-1)^r(-1)^r=1$ (\lean{h1}). Positivity of the ascending symbol
at $\sigma>0$ is \lean{ascPochhammer_pos}, and the final \lean{calc} cancels the two signs, exactly as in the last
display of the human proof. The optimality theorem is one line: the first projection \lean{.1} of the upstream theorem.

\begin{leanbox}[(negative control, \lean{Controls.lean}, verbatim)]
\begin{lstlisting}[language=Lean]
theorem riesz_neg_two_not_admissible : ¬ AdmissiblePotential (TriangularRiesz.riesz (-2)) := by
  intro h
  have h3 := h.2.2 1 1 one_pos
  have hfun : TriangularRiesz.riesz (-2) = fun t => t := by
    funext t
    unfold TriangularRiesz.riesz
    norm_num
  rw [hfun] at h3
  simp at h3 <;> linarith
\end{lstlisting}
\end{leanbox}
The hypothesis $s>0$ cannot be dropped: for $s=-2$ the ``potential'' is $g(t)=t$, whose first derivative is $+1$, so
the condition $(-1)^1g'(1)\ge0$ fails. The control picks exactly that instance (\lean{h.2.2 1 1}: $r=1$, $t=1$). It
shows that \lean{AdmissiblePotential} is not vacuous \tierA.

% =====================================================================================================
\section{Result D1b: the Riesz lattice energy is finite for $s>2$}\label{ch2:sec:d1b}

\paragraph{Informal statement.} For $s>2$ the Epstein-type sum $\sum_{a\in A\setminus\{0\}}\|a\|^{-s}$ converges.

\begin{lemma}[Norm identity]\label{ch2:lem:norm}
$\|p(j,k)\|^2=\dfrac1b\,(j^2+jk+k^2)$ for all $(j,k)\in\Z^2$.
\end{lemma}

\begin{humanproof}
$\|p(j,k)\|^2=b^{-1}\bigl((j+k/2)^2+k^2b^2\bigr)=b^{-1}\bigl(j^2+jk+\tfrac14k^2+\tfrac34k^2\bigr)=b^{-1}(j^2+jk+k^2)$,
using $b^2=3/4$.
\end{humanproof}

\begin{theorem}[D1b]\label{ch2:thm:d1b}
For every $s>2$, $\mathcal{E}_{g_s}(A)=\sum_{a\in A\setminus\{0\}}\|a\|^{-s}<\infty$.
\end{theorem}

\begin{humanproof}
Let $\sigma=s/2>1$ and $Q(j,k)=j^2+jk+k^2$.

\emph{Step 1 (summability over $\Z^2$).} Since $jk\ge-\tfrac12(j^2+k^2)$, we have $Q(j,k)\ge\tfrac12(j^2+k^2)$, so $Q$
vanishes only at $(0,0)$ and $Q^{-\sigma}\le2^\sigma(j^2+k^2)^{-\sigma}$ elsewhere. The $8n$ integer points with
$\max(|j|,|k|)=n\ge1$ satisfy $j^2+k^2\ge n^2$, so
\[
\sum_{(j,k)\neq(0,0)}Q(j,k)^{-\sigma}\le2^\sigma\sum_{n\ge1}8n\cdot n^{-2\sigma}<\infty\qquad(2\sigma-1>1).
\]

\emph{Step 2 (scaling).} By Lemma~\ref{ch2:lem:norm}, $\|p(j,k)\|^{-s}=(Q/b)^{-\sigma}=b^{\sigma}Q^{-\sigma}$, a constant
multiple of a summable family.

\emph{Step 3 (reindexing).} Each $a\in A\setminus\{0\}$ is $p(j,k)$ for some index; choose one, $\iota(a)$. Then
$p(\iota(a))=a$, so $\iota$ is injective (if $\iota(a)=\iota(a')$ then $a=p(\iota(a))=p(\iota(a'))=a'$). A nonnegative
summable family stays summable when restricted along an injective map, so $\sum_{a}\|a\|^{-s}=\sum_a
b^\sigma Q(\iota(a))^{-\sigma}<\infty$.
\end{humanproof}

\begin{leanbox}[(\lean{formal_cert/Triangular/RieszFinite.lean}, verbatim, shortened)]
\begin{lstlisting}[language=Lean,escapeinside={(*}{*)}]
lemma norm_sq_triangularPoint (j k : ℤ) :
    ‖triangularPoint (j, k)‖ ^ 2 = (1 / b) * ((j : ℝ) ^ 2 + (j : ℝ) * k + (k : ℝ) ^ 2) := by
  rw [EuclideanSpace.real_norm_sq_eq]
  -- ...
  field_simp
  rw [hsb, hb2]
  ring

/-- A chosen index of a lattice point. -/
noncomputable def toIdx (a : {x : Plane // x ∈ A ∧ x ≠ 0}) : ℤ × ℤ := Classical.choose a.2.1

lemma toIdx_spec (a : {x : Plane // x ∈ A ∧ x ≠ 0}) : triangularPoint (toIdx a) = a.val :=
  Classical.choose_spec a.2.1

lemma toIdx_injective : Function.Injective toIdx := by
  intro a a' h
  apply Subtype.ext
  rw [(*$\leftarrow$*) toIdx_spec a, (*$\leftarrow$*) toIdx_spec a', h]

/-- **D1b.** For every `s > 2` the Riesz `s`-energy of the density-one triangular lattice is finite. -/
theorem latticeEnergy_riesz_lt_top {s : ℝ} (hs : 2 < s) : latticeEnergy (riesz s) < ⊤ := by
  have hσ : 1 < s / 2 := by linarith
  set Q : ℤ × ℤ → ℝ := fun jk => (jk.1 : ℝ) ^ 2 + (jk.1 : ℝ) * jk.2 + (jk.2 : ℝ) ^ 2 with hQ
  -- ...
  -- (*$\Sigma$*) Q^(-σ) over ℤ² from openai/math, via (j,k) (*$\mapsto$*) (j,-k)
  have hQsum : Summable (fun jk : ℤ × ℤ => Q jk ^ (-(s / 2))) := by
    have hn := LatticeSummability.summable_normForm_neg_rpow hσ
    -- ...
  have hh : Summable (fun jk : ℤ × ℤ => ((1 / b) * Q jk) ^ (-(s / 2))) := by
    -- ...
    exact hQsum.mul_left _
  -- ...
  have hfsum : Summable f := by
    have := hh.comp_injective toIdx_injective
    -- ...
  rw [(*$\leftarrow$*) ENNReal.ofReal_tsum_of_nonneg hf0 hfsum]
  exact ENNReal.ofReal_lt_top
\end{lstlisting}
\end{leanbox}
(The elided lines include an equivalence of $\Z\times\Z$ written with a Unicode arrow that this book's listing style
cannot typeset; read the file for the full text.)

\paragraph{How the Lean proof mirrors the human proof.} The three steps appear in the same order. Step~1 is not reproved.
The openai/math library already contains the lemma
\begin{center}\lean{LatticeSummability.summable_normForm_neg_rpow},\end{center}
\noindent summability of
$(x^2-xy+y^2)^{-\sigma}$ for $\sigma>1$, and the file transports it to $Q$ by the substitution $(j,k)\mapsto(j,-k)$,
which turns $j^2-jk+k^2$ into $j^2+jk+k^2$. (In Lean's real powers $0^{-\sigma}=0$, so the term at $(0,0)$ does no harm
and the sum may run over all of $\Z^2$.) Step~2 is \lean{hh}, a constant multiple via \lean{Real.mul_rpow}. Step~3 is
\lean{toIdx}: an index chosen with \lean{Classical.choose}, whose injectivity is proved exactly as in the human proof,
from $p\circ\iota=\mathrm{id}$, without ever proving that $p$ itself is injective. \lean{Summable.comp_injective}
restricts the summable family, and \lean{ENNReal.ofReal_tsum_of_nonneg} moves the real sum into $[0,\infty]$, where a
real number is $<\infty$.

% =====================================================================================================
\section{Result D2: every density}\label{ch2:sec:d2}

\paragraph{Informal statement.} At every density $\rho>0$ the triangular lattice rescaled to density $\rho$,
$\rho^{-1/2}A$, is optimal for every admissible potential, and it attains the bound.

\begin{theorem}[D2]\label{ch2:thm:d2}
Let $\rho>0$, $h$ admissible and $C$ locally finite of density $\rho$. Then
\[
\sum_{a\in A\setminus\{0\}}h\Bigl(\frac{\|a\|^2}{\rho}\Bigr)\le E_h(C),
\qquad\text{and}\qquad E_h(\rho^{-1/2}A)=\sum_{a\in A\setminus\{0\}}h\Bigl(\frac{\|a\|^2}{\rho}\Bigr).
\]
Consequently $\rho^{-1/2}A$ is locally finite, has density $\rho$, and $E_h(\rho^{-1/2}A)\le E_h(C)$.
\end{theorem}

The proof consists of four scaling lemmas. Throughout, $c>0$ and $cC=\{cx:x\in C\}$.

\begin{lemma}[Scaling counts]\label{ch2:lem:count}
If $C$ is locally finite, so is $cC$, and $N_R(cC)=N_{R/c}(C)$.
\end{lemma}
\begin{humanproof}
$cx\in B_R\iff\|x\|\le R/c$, so $cC\cap B_R=c\,(C\cap B_{R/c})$, the image of a finite set under the injective map
$x\mapsto cx$. It is finite and has the same number of elements.
\end{humanproof}

\begin{lemma}[Scaling density]\label{ch2:lem:dens}
If $C$ is locally finite with density $\rho$, then $\sqrt\rho\,C$ has density one.
\end{lemma}
\begin{humanproof}
With $c=\sqrt\rho$ and Lemma~\ref{ch2:lem:count}, for $R>0$
\[
\frac{N_R(cC)}{\pi R^2}=\frac{N_{R/c}(C)}{\pi R^2}=\frac1\rho\cdot\frac{N_{R/c}(C)}{\pi(R/c)^2}
\xrightarrow[R\to\infty]{}\frac1\rho\cdot\rho=1,
\]
because $R/c\to\infty$ when $R\to\infty$.
\end{humanproof}

\begin{lemma}[Scaling energies]\label{ch2:lem:en}
If $C$ is locally finite, $E_g(cC;R)=E_{g(c^2\,\cdot)}(C;R/c)$ for every $R$, and $E_g(cC)=E_{g(c^2\,\cdot)}(C)$.
\end{lemma}
\begin{humanproof}
The disk sets correspond under $x\mapsto cx$ (Lemma~\ref{ch2:lem:count}), the counts agree, and
$g(\|cx-cy\|^2)=g(c^2\|x-y\|^2)$, which proves the first identity term by term. For the second, put
$F(R)=E_{g(c^2\cdot)}(C;R)$. The map $R\mapsto R/c$ is an increasing bijection of $\R$, so it maps the filter of
neighbourhoods of $+\infty$ onto itself, and $\liminf_{R\to\infty}F(R/c)=\liminf_{R\to\infty}F(R)$.
\end{humanproof}

\begin{lemma}[Rescaled potentials]\label{ch2:lem:adm}
If $g$ is admissible and $c>0$, then $t\mapsto g(ct)$ is admissible.
\end{lemma}
\begin{humanproof}
Smoothness and positivity are preserved because $t>0\Rightarrow ct>0$. By the chain rule and induction on $r$,
$\frac{d^r}{dt^r}g(ct)=c^r g^{(r)}(ct)$, so $(-1)^r\frac{d^r}{dt^r}g(ct)=c^r\cdot(-1)^rg^{(r)}(ct)\ge0$.
\end{humanproof}

\begin{humanproof}[Human proof of Theorem~\ref{ch2:thm:d2}]
\emph{Bound.} Let $c=\sqrt\rho$ and $h_\rho(t)=h(t/\rho)$, admissible by Lemma~\ref{ch2:lem:adm}. The configuration
$cC$ is locally finite with density one (Lemmas~\ref{ch2:lem:count} and~\ref{ch2:lem:dens}). Theorem~\ref{ch2:thm:upstream}
gives $\mathcal{E}_{h_\rho}(A)\le E_{h_\rho}(cC)$. By Lemma~\ref{ch2:lem:en},
$E_{h_\rho}(cC)=E_{h_\rho(c^2\,\cdot)}(C)$, and $h_\rho(c^2t)=h(\rho t/\rho)=h(t)$. So $\mathcal{E}_{h_\rho}(A)\le E_h(C)$,
which is the bound.

\emph{Attainment.} By Lemma~\ref{ch2:lem:en} with $c=\rho^{-1/2}$, $E_h(\rho^{-1/2}A)=E_{h(\rho^{-1}\cdot)}(A)=E_{h_\rho}(A)$,
and the second conjunct of Theorem~\ref{ch2:thm:upstream} for $h_\rho$ gives $E_{h_\rho}(A)=\mathcal{E}_{h_\rho}(A)$.

\emph{Density of $\rho^{-1/2}A$.} With $c=\rho^{-1/2}$, $N_R(cA)/(\pi R^2)=\rho\cdot N_{R\sqrt\rho}(A)/(\pi(R\sqrt\rho)^2)\to\rho$.
\end{humanproof}

\begin{leanbox}[(\lean{formal_cert/Triangular/AnyDensity.lean}, verbatim, shortened)]
\begin{lstlisting}[language=Lean,escapeinside={(*}{*)}]
/-- `C` has density `ρ`. -/
def DensityRho (ρ : ℝ) (C : Set Plane) : Prop :=
  Tendsto (fun R : ℝ => (diskCount C R : ℝ) / (Real.pi * R ^ 2)) atTop (𝓝 ρ)

lemma admissible_comp_mul {g : ℝ → ℝ} (hg : AdmissiblePotential g) {c : ℝ} (hc : 0 < c) :
    AdmissiblePotential (fun t => g (c * t)) := by
  -- ...
lemma diskCount_smul {c : ℝ} (hc : 0 < c) {C : Set Plane} (hC : LocallyFinite C) (R : ℝ) :
    diskCount (c • C) R = diskCount C (R / c) := by
  -- ...
lemma density_smul {ρ : ℝ} (hρ : 0 < ρ) {C : Set Plane} (hC : LocallyFinite C) (hd : DensityRho ρ C) :
    DensityOne (Real.sqrt ρ • C) := by
  -- ...
lemma diskEnergy_smul {c : ℝ} (hc : 0 < c) (g : ℝ → ℝ) {C : Set Plane} (hC : LocallyFinite C) (R : ℝ) :
    diskEnergy g (c • C) R = diskEnergy (fun t => g (c ^ 2 * t)) C (R / c) := by
  -- ... (liminf_comp_div is quoted in full below)
lemma energy_smul {c : ℝ} (hc : 0 < c) (g : ℝ → ℝ) {C : Set Plane} (hC : LocallyFinite C) :
    energy g (c • C) = energy (fun t => g (c ^ 2 * t)) C := by
  -- ...

/-- **D2.** Universal optimality at every density `ρ > 0`. -/
theorem universal_any_density (h : ℝ → ℝ) (hh : AdmissiblePotential h) {ρ : ℝ} (hρ : 0 < ρ) (C : Set Plane)
    (hC : LocallyFinite C) (hd : DensityRho ρ C) :
    latticeEnergy (fun t => h (ρ⁻¹ * t)) ≤ energy h C := by
  have hs : 0 < Real.sqrt ρ := Real.sqrt_pos.2 hρ
  have hg := admissible_comp_mul hh (inv_pos.2 hρ)
  have key := (universal_energy_minimum (fun t => h (ρ⁻¹ * t)) (Real.sqrt ρ • C) hg (locallyFinite_smul hs hC)
    (density_smul hρ hC hd)).1
  rw [energy_smul hs _ hC] at key
  have e : (fun t => (fun t => h (ρ⁻¹ * t)) (Real.sqrt ρ ^ 2 * t)) = h := by
    funext t
    simp only
    rw [Real.sq_sqrt hρ.le, (*$\leftarrow$*) mul_assoc, inv_mul_cancel₀ hρ.ne', one_mul]
  rw [e] at key
  exact key

theorem attained (h : ℝ → ℝ) (hh : AdmissiblePotential h) {ρ : ℝ} (hρ : 0 < ρ) :
    energy h ((Real.sqrt ρ)⁻¹ • A) = latticeEnergy (fun t => h (ρ⁻¹ * t)) := by
  have hs : 0 < Real.sqrt ρ := Real.sqrt_pos.2 hρ
  rw [energy_smul (inv_pos.2 hs) h A_lf, inv_pow, Real.sq_sqrt hρ.le]
  exact ((universal_energy_minimum _ A (admissible_comp_mul hh (inv_pos.2 hρ)) A_lf A_d1).2).symm

/-- **D2, optimality form.** At every density `ρ > 0` the rescaled lattice `ρ^{-1/2} A` has density `ρ`, and its energy is
`≤` that of every locally finite density-`ρ` configuration. -/
theorem triangular_optimal_any_density (h : ℝ → ℝ) (hh : AdmissiblePotential h) {ρ : ℝ} (hρ : 0 < ρ) :
    LocallyFinite ((Real.sqrt ρ)⁻¹ • A) ∧ DensityRho ρ ((Real.sqrt ρ)⁻¹ • A) ∧
      ∀ C : Set Plane, LocallyFinite C → DensityRho ρ C → energy h ((Real.sqrt ρ)⁻¹ • A) ≤ energy h C :=
  ⟨locallyFinite_smul (inv_pos.2 (Real.sqrt_pos.2 hρ)) A_lf, scaled_density hρ, fun C hC hd => by
    rw [attained h hh hρ]; exact universal_any_density h hh hρ C hC hd⟩
\end{lstlisting}
The reparametrisation of the $\liminf$ (Lemma~\ref{ch2:lem:en}, second part) is the following, verbatim:
\begin{lstlisting}[language=Lean,escapeinside={(*}{*)}]
lemma liminf_comp_div {c : ℝ} (hc : 0 < c) (F : ℝ → ENNReal) :
    liminf (fun R => F (R / c)) atTop = liminf F atTop := by
  have hmap : map (fun R : ℝ => R / c) atTop = atTop := by
    have h1 := (OrderIso.mulLeft₀ c⁻¹ (inv_pos.2 hc)).map_atTop
    have h2 : ((OrderIso.mulLeft₀ c⁻¹ (inv_pos.2 hc)) : ℝ → ℝ) = fun R => R / c := by
      funext R
      exact (div_eq_inv_mul R c).symm
    rw [(*$\leftarrow$*) h2]
    exact h1
  rw [show (fun R => F (R / c)) = F ∘ (fun R : ℝ => R / c) from rfl]
  unfold Filter.liminf
  rw [(*$\leftarrow$*) Filter.map_map, hmap]
\end{lstlisting}
\end{leanbox}

\paragraph{How the Lean proof mirrors the human proof.} The four lemmas appear under their own names:
\lean{diskCount_smul} (counts), \lean{density_smul} (density), \lean{diskEnergy_smul} together with
\lean{liminf_comp_div} and \lean{energy_smul} (energies), and \lean{admissible_comp_mul} (potentials; its proof is the
induction $\frac{d^r}{dt^r}g(ct)=c^rg^{(r)}(ct)$, stated as \lean{iterate_deriv_comp_mul}). The main proof is the
human ``Bound'' paragraph line by line: \lean{key} is the upstream bound for $h_\rho$ and $\sqrt\rho\,C$;
\lean{energy_smul} rewrites the energy of $\sqrt\rho\,C$; \lean{e} is the computation $h_\rho(c^2t)=h(t)$. The
$\liminf$ lemma shows a typical formalisation move: ``an increasing bijection maps $+\infty$ to $+\infty$'' becomes
\lean{OrderIso.map_atTop}, and the $\liminf$ of a composition becomes the $\liminf$ along the image filter
(\lean{Filter.map_map}).

\subsection{Why the factor $\rho^{-1}$ is needed: a negative control}

It is tempting to state D2 as ``the density-one lattice energy $\mathcal{E}_h(A)$ is a lower bound at every density''.
That is false. At density $\rho<1$ the points are farther apart, so repulsive energies are smaller.

\begin{proposition}[\lean{unscaled_false}]\label{ch2:prop:unscaled}
It is not true that $\mathcal{E}_{g_4}(A)\le E_{g_4}(C)$ for every locally finite $C$ of density $1/4$.
\end{proposition}
\begin{humanproof}
Take $C=2A=(1/4)^{-1/2}A$, which has density $1/4$ by Theorem~\ref{ch2:thm:d2}. By attainment,
$E_{g_4}(2A)=\sum_{a\neq0}(4\|a\|^2)^{-2}=\tfrac1{16}\sum_{a\ne0}\|a\|^{-4}$. The sum is finite (D1b with $s=4$) and
positive, so $E_{g_4}(2A)<\mathcal{E}_{g_4}(A)$.
\end{humanproof}

In Lean the strict inequality is proved termwise rather than by the factor $1/16$: every term
$((1/4)^{-1}\|a\|^2)^{-2}$ is $\le$, and the term at $p(1,0)$ is $<$, the corresponding term $(\|a\|^2)^{-2}$; since the
right-hand sum is finite (\lean{latticeEnergy_riesz_lt_top}), \lean{ENNReal.tsum_lt_tsum} gives a strict inequality of
the sums.
\begin{leanbox}[(\lean{AnyDensity.lean}, verbatim, shortened)]
\begin{lstlisting}[language=Lean,escapeinside={(*}{*)}]
theorem unscaled_false :
    ¬ (∀ C : Set Plane, LocallyFinite C → DensityRho (1 / 4) C →
        latticeEnergy (TriangularRiesz.riesz 4) ≤ energy (TriangularRiesz.riesz 4) C) := by
  intro H
  have hρ : (0 : ℝ) < 1 / 4 := by norm_num
  have key := H _ (locallyFinite_smul (inv_pos.2 (Real.sqrt_pos.2 hρ)) A_lf) (scaled_density hρ)
  rw [attained _ (TriangularRiesz.riesz_admissible (by norm_num)) hρ] at key
  -- key : (*$\Sigma$*) ofReal(t^{-2}) ≤ (*$\Sigma$*) ofReal((4t)^{-2}); but (4t)^{-2} < t^{-2} termwise and the sum is finite
  have hfin := TriangularRiesz.latticeEnergy_riesz_lt_top (s := 4) (by norm_num)
  -- ...
  have strict := ENNReal.tsum_lt_tsum hfin2 hle (hlt ⟨_, hpt⟩)
  exact absurd key (not_le.2 strict)
\end{lstlisting}
\end{leanbox}

For the Riesz family the correct bound has a closed form. Since $g_s(t/\rho)=\rho^{s/2}g_s(t)$,
\[
\mathcal{E}_{g_s(\rho^{-1}\cdot)}(A)=\rho^{s/2}\,\mathcal{E}_{g_s}(A),
\qquad\text{so}\qquad \rho^{s/2}\sum_{a\in A\setminus\{0\}}\|a\|^{-s}\le E_{g_s}(C)
\]
for every locally finite $C$ of density $\rho$ \tierA. In Lean these are \lean{riesz_lattice_homogeneous} and
\lean{riesz_any_density_homogeneous} in \lean{Controls.lean}. At $\rho=1/4$, $s=4$ the factor is $1/16$, as in the proof above. The control
\lean{rho_one} checks that at $\rho=1$ D2 gives back the density-one statement, and \lean{densityRho_one_iff} that
\lean{DensityRho 1} \emph{is} \lean{DensityOne} (by \lean{Iff.rfl}) \tierA.

% =====================================================================================================
\section{Result D3: universal optimality, and a second minimizer}\label{ch2:sec:d3}

\subsection{The definition}
\begin{definition}[Universally optimal, lower-energy form]\label{ch2:def:uo}
A configuration $C$ is \emph{universally optimal at density $\rho$} if $C$ is locally finite, has density $\rho$, and
$E_g(C)\le E_g(C')$ for every admissible $g$ and every locally finite $C'$ of density $\rho$.
\end{definition}

This transcribes Definition~1.3 of~\cite{ckmrv} into upstream's notions of density and lower energy. It is
\emph{weaker} than CKMRV's notion: CKMRV also require the energy of $C$ itself to exist as a limit, while here only the
$\liminf$ enters \tierC. A Lean theorem about \lean{UniversallyOptimal} therefore says nothing about whether the limit
exists.

\begin{theorem}[\lean{triangular_universally_optimal}]\label{ch2:thm:uo}
For every $\rho>0$, $\rho^{-1/2}A$ is universally optimal at density $\rho$.
\end{theorem}
\begin{humanproof}
This is the last sentence of Theorem~\ref{ch2:thm:d2}, read against Definition~\ref{ch2:def:uo}.
\end{humanproof}

\begin{leanbox}[(\lean{formal_cert/Triangular/UniversalOptimality.lean}, verbatim)]
\begin{lstlisting}[language=Lean]
/-- CKMRV Definition 1.3, in upstream's density and lower-energy notions. -/
def UniversallyOptimal (ρ : ℝ) (C : Set Plane) : Prop :=
  LocallyFinite C ∧ DensityRho ρ C ∧
    ∀ g : ℝ → ℝ, AdmissiblePotential g → ∀ C' : Set Plane, LocallyFinite C' → DensityRho ρ C' →
      energy g C ≤ energy g C'

/-- **Universal optimality of the triangular lattice at every density** (conditional on upstream). -/
theorem triangular_universally_optimal {ρ : ℝ} (hρ : 0 < ρ) :
    UniversallyOptimal ρ ((Real.sqrt ρ)⁻¹ • A) := by
  refine ⟨locallyFinite_smul (inv_pos.2 (Real.sqrt_pos.2 hρ)) A_lf, scaled_density hρ, ?_⟩
  intro g hg C' hC' hd
  exact (triangular_optimal_any_density g hg hρ).2.2 C' hC' hd
\end{lstlisting}
\end{leanbox}

\subsection{The vacancy theorem}
Let $V=A\setminus\{0\}$, the lattice with one point removed (a \emph{vacancy}).

\begin{theorem}[Vacancy]\label{ch2:thm:vac}
\begin{enumerate}[label=(\roman*)]
\item $V$ is locally finite and has density one.
\item $E_g(V)=\mathcal{E}_g(A)$ for every admissible $g$.
\item There is no isometry $f$ of the plane with $f(A)=V$.
\end{enumerate}
Hence $V$ is universally optimal at density one but is not an isometric copy of $A$: universal optima, in the
lower-energy sense, are not unique up to isometry among all locally finite configurations.
\end{theorem}

\begin{humanproof}
Write $N=N_R(A)$, and assume $R\ge0$, so that $0\in A\cap B_R$ and $N\ge1$.

\emph{(i) Counts and density.} $V\subseteq A$, so $V$ is locally finite, and $V\cap B_R=(A\cap B_R)\setminus\{0\}$ has
$N-1$ elements. Hence
\[
\frac{N_R(V)}{\pi R^2}=\frac{N_R(A)}{\pi R^2}-\frac1{\pi R^2}\longrightarrow1-0=1.
\]

\emph{(ii), upper bound.} Denote by $S_V(R)$ and $S_A(R)$ the two double sums in the disk energies. Every ordered pair
$(x,y)$ of distinct points of $V\cap B_R$ is also a pair of $A\cap B_R$, and all terms are $\ge0$, so $S_V(R)\le
S_A(R)$. Since $N\to\infty$ (density one forces $N\sim\pi R^2$), for large $R$ we have $N\ge2$ and
\[
E_g(V;R)=\frac{S_V(R)}{N-1}\le\frac{N}{N-1}\cdot\frac{S_A(R)}{N}=\kappa(R)\,E_g(A;R),\qquad
\kappa(R)=\frac{N}{N-1}\longrightarrow1.
\]
Now we take lower limits. If $L=\liminf_R E_g(A;R)=\infty$ there is nothing to prove. Otherwise choose $R_n\to\infty$
with $E_g(A;R_n)\to L$. Then $\kappa(R_n)E_g(A;R_n)\to1\cdot L$, so
\[
E_g(V)=\liminf_RE_g(V;R)\le\liminf_n\kappa(R_n)E_g(A;R_n)=L=E_g(A)=\mathcal{E}_g(A),
\]
the last equality by the second conjunct of Theorem~\ref{ch2:thm:upstream}.

\emph{(ii), lower bound.} By (i), $V$ is an admissible competitor: locally finite of density one. The first conjunct of
Theorem~\ref{ch2:thm:upstream} gives $\mathcal{E}_g(A)\le E_g(V)$. (Notice that this direction uses the universal
theorem itself; we do not compute $E_g(V)$ from below by hand.)

\emph{(iii) No isometry.} Suppose $f$ is an isometry of $\R^2$ with $f(A)=V$. An isometry of a strictly convex real
normed space, in particular of the Euclidean plane, is affine, so
\[
f(2a-b)=2f(a)-f(b)\qquad\text{for all }a,b.
\]
The points $p(1,0)$ and $p(2,0)$ are nonzero (their squared norms are $1/b$ and $4/b$ by Lemma~\ref{ch2:lem:norm}), so
they lie in $V=f(A)$: $p(1,0)=f(a)$ and $p(2,0)=f(b)$ with $a,b\in A$. The lattice is closed under integer combinations,
so $2a-b\in A$, and
\[
f(2a-b)=2p(1,0)-p(2,0)=p(0,0)=0.
\]
So $0\in f(A)=V$, which contradicts $0\notin V$. In words: $p(1,0)$ is the midpoint of $0$ and $p(2,0)$; an affine
bijection preserves midpoints and the lattice is closed under ``reflect $b$ through $a$'', so the image of $A$ must
contain the point that the vacancy has removed.

\emph{Conclusion.} By (i) and (ii), for every admissible $g$ and every competitor $C'$,
$E_g(V)=\mathcal{E}_g(A)\le E_g(C')$; so $V$ is universally optimal at density one. With (iii) this is the
non-uniqueness statement.
\end{humanproof}

\begin{leanbox}[(\lean{UniversalOptimality.lean}, verbatim, shortened)]
\begin{lstlisting}[language=Lean,escapeinside={(*}{*)}]
def vacancy : Set Plane := A \ {0}

lemma diskCount_vacancy {R : ℝ} (hR : 0 ≤ R) : diskCount vacancy R = diskCount A R - 1 := by
  -- ...
/-- The correction factor `N / (N - 1)`. -/
noncomputable def corr (R : ℝ) : ENNReal :=
  ENNReal.ofReal ((diskCount A R : ℝ) / ((diskCount A R : ℝ) - 1))

lemma corr_tendsto : Tendsto corr atTop (𝓝 1) := by
  -- ...
lemma diskEnergy_vacancy_le (g : ℝ → ℝ) :
    (*$\forall^{\mathsf f}$*) R in atTop, diskEnergy g vacancy R ≤ corr R * diskEnergy g A R := by
  -- ...

/-- **The vacancy configuration is also a minimiser**, for every admissible `g`. -/
theorem vacancy_energy (g : ℝ → ℝ) (hg : AdmissiblePotential g) : energy g vacancy = latticeEnergy g := by
  apply le_antisymm
  · have hA := (universal_energy_minimum g A hg A_lf A_d1).2
    have hc := corr_tendsto
    calc energy g vacancy = liminf (diskEnergy g vacancy) atTop := rfl
      _ ≤ liminf (corr * diskEnergy g A) atTop := liminf_le_liminf (diskEnergy_vacancy_le g)
      _ ≤ limsup corr atTop * liminf (diskEnergy g A) atTop := by
          apply ENNReal.liminf_mul_le <;> left <;> rw [hc.limsup_eq] <;> simp
      _ = latticeEnergy g := by
          rw [hc.limsup_eq, one_mul, hA]
          rfl
  · exact (universal_energy_minimum g vacancy hg vacancy_lf vacancy_density).1

theorem vacancy_not_isometric_image : ¬ ∃ f : Plane → Plane, Isometry f ∧ f '' A = vacancy := by
  rintro ⟨f, hf, hfA⟩
  set φ := hf.affineIsometryOfStrictConvexSpace
  have hφ : ∀ p, φ p = f p := fun p => rfl
  have key : ∀ a b : Plane, f ((2 : ℝ) • a - b) = (2 : ℝ) • f a - f b := by
    -- ...
  have hp : triangularPoint (1, 0) ∈ vacancy := by
    -- ...
  have hq : triangularPoint (2, 0) ∈ vacancy := by
    -- ...
  rw [(*$\leftarrow$*) hfA] at hp hq
  obtain ⟨a, ha, hfa⟩ := hp
  obtain ⟨b, hb, hfb⟩ := hq
  have h0 : f ((2 : ℝ) • a - b) = 0 := by
    rw [key, hfa, hfb, (*$\leftarrow$*) triangularPoint_comb]
    convert triangularPoint_zero using 2 <;> norm_num
  have hmem : (0 : Plane) ∈ f '' A := ⟨_, A_comb ha hb, h0⟩
  rw [hfA] at hmem
  exact hmem.2 rfl

/-- **Non-uniqueness.** There is a universal optimum at density one that is not an isometric copy of the triangular
lattice: the minimiser is not determined, only the minimum value. -/
theorem universal_optimum_not_unique :
    ∃ C : Set Plane, UniversallyOptimal 1 C ∧ ¬ ∃ f : Plane → Plane, Isometry f ∧ f '' A = C :=
  ⟨vacancy, vacancy_universally_optimal, vacancy_not_isometric_image⟩
\end{lstlisting}
(Typographical note: in the file, the quantifier of \lean{diskEnergy_vacancy_le} is Lean's ``eventually'' binder, written
with a superscript letter f after $\forall$; the listing shows it as $\forall^{\mathsf f}$.)
\end{leanbox}

\paragraph{How the Lean proof mirrors the human proof.} \lean{diskCount_vacancy} is the count $N-1$;
\lean{vacancy_density} is the limit $1-0$. The factor $\kappa$ is \lean{corr}, and \lean{corr_tendsto} is
$\kappa\to1$. \lean{diskEnergy_vacancy_le} is the inequality $E_g(V;R)\le\kappa(R)E_g(A;R)$ for all large $R$, built
from \lean{pair_sum_le} ($S_V\le S_A$, by dropping terms) and \lean{coeff_eq} ($\frac1{N-1}=\frac N{N-1}\cdot\frac1N$ in
$[0,\infty]$). The upper bound in \lean{vacancy_energy} is a three-step \lean{calc}. The human proof's subsequence
argument is replaced by the library inequality \lean{ENNReal.liminf_mul_le},
$\liminf(uv)\le\limsup u\cdot\liminf v$ in $[0,\infty]$, whose side conditions ($\limsup u\neq0$ or $\liminf v\neq\infty$,
and the symmetric one) are discharged because $\limsup\kappa=1$. The lower bound is literally ``apply the upstream
theorem to $V$''. In \lean{vacancy_not_isometric_image}, \lean{Isometry.affineIsometryOfStrictConvexSpace} is the
statement ``isometries are affine'', \lean{key} is $f(2a-b)=2f(a)-f(b)$, \lean{triangularPoint_comb} and \lean{A_comb}
say that the lattice is closed under $(a,b)\mapsto2a-b$, and the last three lines derive $0\in V$ and the contradiction.

\subsection{Interpretation and what remains open}
\begin{itemize}
\item \tierC{} The lower energy cannot see a single missing point, and the same argument applies to any finite set of
  missing points (Exercise~\ref{ch2:ex:finite}). So any statement of the form ``the minimizer is the triangular
  lattice'' is \emph{false} among all locally finite configurations. This is the remark after CKMRV Theorem~1.4
  (``removing a single particle changes neither the density nor the energy'')~\cite{ckmrv} \tierL, here in Lean
  \tierA{} (for the lower-energy notion).
\item \tierC{} The vacancy is not periodic. The meaningful uniqueness question is therefore uniqueness among
  \emph{periodic} configurations, the form CKMRV prove in dimensions $8$ and $24$. In dimension~$2$ it is \textbf{open},
  and nothing in this book addresses it.
\item \tierC{} For physics (Chapter~4), the theorem says that the triangular lattice has the least energy per particle
  in the bulk; it does not say that a physical system will be found in that state, or that defects cost energy per
  particle. A single vacancy changes the total pair energy of a large disk by an amount that stays bounded (for
  $s>2$ Riesz potentials, by the argument of Exercise~\ref{ch2:ex:add}), and the energy per particle does not see it.
\end{itemize}

% =====================================================================================================
\section{Result D4: open versus closed disks}\label{ch2:sec:d4}

Upstream measures densities and energies on \emph{closed} disks $B_R$. A reader might worry that open disks
$B_R^\circ=\{x:\|x\|<R\}$ give a different theory: on the boundary circle a lattice can carry many points. We now prove
that, for locally finite configurations, nothing changes.

\paragraph{Notation.} For locally finite $C$ write $N^\circ_R(C)=\#(C\cap B^\circ_R)$ and define $E^\circ_g(C;R)$ and
$E^\circ_g(C)=\liminf_RE^\circ_g(C;R)$ exactly as in Definition~\ref{ch2:def:energy}, with $B^\circ_R$ in place of $B_R$.
$C$ has \emph{open density} $\rho$ if $N^\circ_R(C)/(\pi R^2)\to\rho$.

\paragraph{Informal statement.} For a locally finite configuration, open and closed disks give the same density and the
same lower energy. Hence upstream's theorem holds verbatim with open disks.

\begin{theorem}[D4]\label{ch2:thm:d4}
Let $C$ be locally finite, $g:\R\to\R$ any function and $\rho\in\R$. Then
\begin{enumerate}[label=(\roman*)]
\item $C$ has density $\rho$ if and only if it has open density $\rho$;
\item $E^\circ_g(C)=E_g(C)$.
\end{enumerate}
\end{theorem}

Lean covers (i) for $\rho=1$ only (\lean{densityOne_iff_open}) and (ii) (\lean{energy_eq_openEnergy});
\lean{OpenDiskAny.lean} proves both without the local-finiteness hypothesis, and \lean{ctl_open} and \lean{ctl_closed} show
the identity is not definitional (a point on the unit circle). Case $\rho\ne1$: \tierC{} (the same argument applies,
not formalised).

Status: \tierC{} for the human proof below. Lean status of D4: \DFourStatus.

The disk energy at radius $R$ depends on $R$ only through the finite set $C\cap B_R$ (respectively $C\cap B^\circ_R$):
the count is its cardinality and the double sum runs over it. The whole proof rests on comparing these sets.

\begin{lemma}[Comparing point sets]\label{ch2:lem:sets}
Let $C$ be locally finite and $R\in\R$. There are radii $\varphi(R)\in[R-1,R]$ and $\psi(R)\in[R,R+1]$ with
\[
C\cap B^\circ_R=C\cap B_{\varphi(R)}\qquad\text{and}\qquad C\cap B_R=C\cap B^\circ_{\psi(R)}.
\]
\end{lemma}

\begin{humanproof}
\emph{First identity.} The set $F=C\cap B^\circ_R$ is finite, since it lies in $C\cap B_R$. Let $m=\max\{\|x\|:x\in F\}$
if $F\neq\emptyset$, and $m=-\infty$ otherwise. Every $x\in F$ has $\|x\|<R$, so $m<R$. Put
$\varphi(R)=\max(m,R-\tfrac12)$, so $R-1\le\varphi(R)<R$. If $x\in C$ and $\|x\|\le\varphi(R)$, then $\|x\|<R$ and
$x\in F$. Conversely, if $x\in F$ then $\|x\|\le m\le\varphi(R)$. So $F=C\cap B_{\varphi(R)}$.

\emph{Second identity.} The set $G=\{x\in C:R<\|x\|\le R+1\}$ is finite, since it lies in $C\cap B_{R+1}$. Let
$m'=\min\{\|x\|:x\in G\}$ if $G\neq\emptyset$, and $m'=+\infty$ otherwise; then $m'>R$. Put
$\psi(R)=\min(m',R+\tfrac12)$, so $R<\psi(R)\le R+1$. If $\|x\|\le R$ then $\|x\|<\psi(R)$. Conversely, let $x\in C$ with
$\|x\|<\psi(R)\le R+\frac12$. If we had $\|x\|>R$, then $x\in G$ and $\|x\|\ge m'\ge\psi(R)$, a contradiction; so
$\|x\|\le R$. Hence $C\cap B_R=C\cap B^\circ_{\psi(R)}$.
\end{humanproof}

Finiteness is used twice, to make the maximum and the minimum exist. Without it the lemma fails
(Exercise~\ref{ch2:ex:lf}).

\begin{lemma}[Lower limits under reparametrisation]\label{ch2:lem:repar}
Let $u:\R\to[0,\infty]$ and $\varphi:\R\to\R$ with $\varphi(R)\to\infty$ as $R\to\infty$. Then
$\liminf_{R\to\infty}u(\varphi(R))\ge\liminf_{R\to\infty}u(R)$.
\end{lemma}

\begin{humanproof}
Recall $\liminf_Ru(R)=\sup_X\inf_{x\ge X}u(x)$. Fix $X$ and put $m_X=\inf_{x\ge X}u(x)$. Since $\varphi(R)\to\infty$ there
is $Y$ with $\varphi(R)\ge X$ for all $R\ge Y$, hence $u(\varphi(R))\ge m_X$ for all $R\ge Y$. Therefore
$\inf_{R\ge Y}u(\varphi(R))\ge m_X$, and so $\liminf_Ru(\varphi(R))\ge m_X$. Taking the supremum over $X$ gives the claim.
\end{humanproof}

\begin{humanproof}[Human proof of Theorem~\ref{ch2:thm:d4}]
\emph{(ii) Energies.} Put $u(R)=E_g(C;R)$ and $u^\circ(R)=E^\circ_g(C;R)$. Because the disk energy depends only on the
point set, Lemma~\ref{ch2:lem:sets} gives
\[
u^\circ(R)=u(\varphi(R)),\qquad u(R)=u^\circ(\psi(R))\qquad\text{for every }R.
\]
Since $\varphi(R)\ge R-1$ and $\psi(R)\ge R$, both reparametrisations tend to $\infty$. Lemma~\ref{ch2:lem:repar} applied
twice gives
\begin{gather*}
E^\circ_g(C)=\liminf_{R\to\infty}u(\varphi(R))\ \ge\ \liminf_{R\to\infty}u(R)=E_g(C),\\
E_g(C)=\liminf_{R\to\infty}u^\circ(\psi(R))\ \ge\ \liminf_{R\to\infty}u^\circ(R)=E^\circ_g(C).
\end{gather*}
Hence $E^\circ_g(C)=E_g(C)$.

\emph{(i) Densities.} Since $B_{R-1}\subseteq B^\circ_R\subseteq B_R\subseteq B^\circ_{R+1}$,
\[
N_{R-1}(C)\le N^\circ_R(C)\le N_R(C)\le N^\circ_{R+1}(C)\qquad(R\ge1).
\]
Suppose $C$ has density $\rho$. For $R>1$,
\[
\frac{N_{R-1}(C)}{\pi R^2}=\frac{N_{R-1}(C)}{\pi(R-1)^2}\cdot\frac{(R-1)^2}{R^2}\longrightarrow\rho\cdot1=\rho,
\]
and $N_R(C)/(\pi R^2)\to\rho$. The middle term $N^\circ_R(C)/(\pi R^2)$ is squeezed between them, so $C$ has open density
$\rho$. Conversely, suppose $C$ has open density $\rho$. Then $N^\circ_R(C)\le N_R(C)\le N^\circ_{R+1}(C)$ and
$N^\circ_{R+1}(C)/(\pi R^2)=\frac{N^\circ_{R+1}(C)}{\pi(R+1)^2}\cdot\frac{(R+1)^2}{R^2}\to\rho$; the same squeeze gives
density $\rho$.
\end{humanproof}

\begin{corollary}[Upstream with open disks]\label{ch2:cor:open}
Let $g$ be admissible and $C$ locally finite with open density one. Then $\mathcal{E}_g(A)\le E^\circ_g(C)$ and
$\mathcal{E}_g(A)=E^\circ_g(A)$.
\end{corollary}
\begin{humanproof}
By Theorem~\ref{ch2:thm:d4}(i), $C$ has density one, so Theorem~\ref{ch2:thm:upstream} gives
$\mathcal{E}_g(A)\le E_g(C)=E^\circ_g(C)$, the equality by Theorem~\ref{ch2:thm:d4}(ii). The lattice $A$ is locally finite,
so $\mathcal{E}_g(A)=E_g(A)=E^\circ_g(A)$.
\end{humanproof}

Status: \tierC{} for the corollary and its proof. Lean status of D4: \DFourStatus.

\begin{remark}
The preregistration of D4 (\lean{results/discovery/D4/preregistration.json}) was committed before the Lean file was
written. It names the hypothesis, the kill criterion (``the liminf comparison for ENNReal-valued step functions on atTop
(real index) cannot be closed in-session''), the controls (axiom audit with planted \lean{sorry}/axiom; non-vacuity at
$C=A$) and a prior probability of about $65\%$ that the Lean proof would go through. The proof went through: the
statements below were kernel-checked, with the footprints recorded in \S\ref{ch2:sec:audit}.
\end{remark}

\begin{leanbox}[(\lean{formal_cert/Triangular/OpenDisk.lean}, verbatim, proofs elided)]
\begin{lstlisting}[language=Lean]
/-- **D4, energy.** The lower energy is the same with open and with closed disks. -/
theorem energy_eq_openEnergy (g : ℝ → ℝ) {C : Set Plane} (hC : LocallyFinite C) :
    energy g C = openEnergy g C := by
  -- ...
/-- **D4, density.** Centred density one is the same with open and with closed disks. -/
theorem densityOne_iff_open {C : Set Plane} (hC : LocallyFinite C) : DensityOne C ↔ DensityOneOpen C := by
  -- ...
/-- **D4.** Upstream's universal optimality theorem, verbatim with open disks. -/
theorem universal_energy_minimum_open (g : ℝ → ℝ) (C : Set Plane) (hg : AdmissiblePotential g)
    (hC : LocallyFinite C) (hd : DensityOneOpen C) :
    latticeEnergy g ≤ openEnergy g C ∧ latticeEnergy g = openEnergy g A := by
  -- ...
theorem A_densityOneOpen : DensityOneOpen A :=
  -- ...
\end{lstlisting}
\end{leanbox}

\begin{leanbox}[(\lean{formal_cert/Triangular/OpenDiskAny.lean}, verbatim, proofs elided)]
\begin{lstlisting}[language=Lean]
theorem energy_eq_openEnergy_any (g : ℝ → ℝ) (C : Set Plane) : energy g C = openEnergy g C := by
  -- ...
theorem densityOne_iff_open_any (C : Set Plane) : DensityOne C ↔ DensityOneOpen C := by
  -- ...
lemma ctl_open : openDiskCount {e} 1 = 0 := by
  -- ...
lemma ctl_closed : diskCount {e} 1 = 1 := by
  -- ...
\end{lstlisting}
\end{leanbox}
Here \lean{e} is \lean{EuclideanSpace.single 0 1}, a point of norm one.

% =====================================================================================================
\section{The axiom audit and the controls}\label{ch2:sec:audit}

A Lean file that ``compiles'' has not thereby been proved. Two failure modes exit with code~$0$:
\begin{itemize}
\item a proof containing \lean{sorry} compiles with a \emph{warning}, not an error;
\item a declared \lean{axiom} can prove anything, including \lean{False}, and compiles silently.
\end{itemize}
Our gate therefore reads the axiom footprint of every theorem with \lean{#print axioms} and accepts a declaration only if
its footprint is contained in the whitelist \{\lean{propext}, \lean{Classical.choice}, \lean{Quot.sound}\}; the
presence of \lean{sorryAx} is a rejection. To show that the gate can fail, the audit file plants one \lean{sorry} and
one false axiom, and the gate must flag both.

\begin{leanbox}[(\lean{formal_cert/Triangular/Axioms.lean}, verbatim, shortened)]
\begin{lstlisting}[language=Lean]
import Triangular.Riesz
import Triangular.RieszFinite
import Triangular.AnyDensity
import Triangular.Controls
import Triangular.UniversalOptimality
import Triangular.OpenDisk
import Triangular.OpenDiskAny

/-! Axiom audit of the D1 corollary, with planted controls. -/

#print axioms TriangularRiesz.riesz_admissible
#print axioms TriangularRiesz.triangular_riesz_optimal
#print axioms OAI.AtomicTriangular.universal_energy_minimum
-- ... (22 audited declarations for D1-D3; 30 in total)
#print axioms TriangularUniversalOptimality.universal_optimum_not_unique
#print axioms TriangularOpenDisk.energy_eq_openEnergy
#print axioms TriangularOpenDisk.densityOne_iff_open
#print axioms TriangularOpenDisk.universal_energy_minimum_open
#print axioms TriangularOpenDisk.A_densityOneOpen
#print axioms TriangularOpenDiskAny.energy_eq_openEnergy_any
#print axioms TriangularOpenDiskAny.densityOne_iff_open_any
#print axioms TriangularOpenDiskAny.ctl_open
#print axioms TriangularOpenDiskAny.ctl_closed

theorem ctl_sorry : TriangularRiesz.riesz 1 1 = 0 := by sorry
axiom ctl_cheat : False
theorem ctl_axiom : TriangularRiesz.riesz 1 1 = 0 := ctl_cheat.elim
#print axioms ctl_sorry
#print axioms ctl_axiom
\end{lstlisting}
\end{leanbox}

Note that the two control theorems state something \emph{false}: $g_1(1)=1^{-1/2}=1\neq0$. Both nevertheless compile.
This is the point of the controls.

\begin{numbox}[(recorded compiler output, not a number but a log)]
The following is the output stored in \lean{results/discovery/D1/compile_Axioms.json} (fields \lean{rc}~$=0$,
\lean{seconds}~$=15.4$) from compiling \lean{Axioms.lean} with the pinned toolchain. We did not rerun the compile for
this book (the machine is shared); the log is quoted, shortened, verbatim.
\begin{lstlisting}
'TriangularRiesz.riesz_admissible' depends on axioms: [propext, Classical.choice, Quot.sound]
'TriangularRiesz.triangular_riesz_optimal' depends on axioms: [propext, Classical.choice, Quot.sound]
'OAI.AtomicTriangular.universal_energy_minimum' depends on axioms: [propext, Classical.choice, Quot.sound]
'TriangularRiesz.norm_sq_triangularPoint' depends on axioms: [propext, Classical.choice, Quot.sound]
'TriangularRiesz.latticeEnergy_riesz_lt_top' depends on axioms: [propext, Classical.choice, Quot.sound]
'OAI.SevenEighths.LatticeSummability.summable_normForm_neg_rpow' depends on axioms: [propext,
 Classical.choice,
 Quot.sound]
'TriangularDensity.universal_any_density' depends on axioms: [propext, Classical.choice, Quot.sound]
...
'TriangularUniversalOptimality.vacancy_not_isometric_image' depends on axioms: [propext, Classical.choice, Quot.sound]
'TriangularUniversalOptimality.vacancy_universally_optimal' depends on axioms: [propext, Classical.choice, Quot.sound]
'TriangularUniversalOptimality.universal_optimum_not_unique' depends on axioms: [propext, Classical.choice, Quot.sound]
'TriangularOpenDisk.energy_eq_openEnergy' depends on axioms: [propext, Classical.choice, Quot.sound]
'TriangularOpenDisk.densityOne_iff_open' depends on axioms: [propext, Classical.choice, Quot.sound]
'TriangularOpenDisk.universal_energy_minimum_open' depends on axioms: [propext, Classical.choice, Quot.sound]
'TriangularOpenDisk.A_densityOneOpen' depends on axioms: [propext, Classical.choice, Quot.sound]
'TriangularOpenDiskAny.energy_eq_openEnergy_any' depends on axioms: [propext, Classical.choice, Quot.sound]
'TriangularOpenDiskAny.densityOne_iff_open_any' depends on axioms: [propext, Classical.choice, Quot.sound]
'TriangularOpenDiskAny.ctl_open' depends on axioms: [propext, Classical.choice, Quot.sound]
'TriangularOpenDiskAny.ctl_closed' depends on axioms: [propext, Classical.choice, Quot.sound]
.../formal_cert/Triangular/Axioms.lean:42:8: warning: declaration uses `sorry`
'ctl_sorry' depends on axioms: [propext, sorryAx, Classical.choice, Quot.sound]
'ctl_axiom' depends on axioms: [ctl_cheat, propext, Classical.choice, Quot.sound]
\end{lstlisting}
\end{numbox}

\paragraph{How to read the log.}
\begin{itemize}
\item The return code is $0$ although the file contains a \lean{sorry} and a false axiom. An exit code is not a proof
  gate.
\item All 30 audited declarations (22 for D1--D3, 8 for D4), including upstream's \lean{universal_energy_minimum}, report the footprint
  [\lean{propext}, \lean{Classical.choice}, \lean{Quot.sound}]: the three standard axioms of Lean's classical
  mathematics, which Mathlib uses everywhere. This is \emph{not} ``zero axioms''. \tierA
\item \lean{ctl_sorry} reports \lean{sorryAx} (and the compiler warns, at line~42), so the gate rejects it.
  \lean{ctl_axiom} reports \lean{ctl_cheat}, outside the whitelist, so the gate rejects it. Both controls are flagged,
  which shows that the gate is able to fail.
\item A clean footprint certifies that the proofs are complete relative to the definitions. It does not certify that
  the definitions mean what the English says (statement adequacy), which is a separate human audit and has not been
  done.
\end{itemize}

\paragraph{The other controls, collected.} Every positive result in this chapter is accompanied by a control showing that
its hypotheses or conclusions are not vacuous \tierA:
\begin{itemize}
\item \lean{riesz_neg_two_not_admissible}: the admissibility hypothesis can fail (\S\ref{ch2:sec:d1});
\item \lean{riesz_sq}: the Lean potential is the Riesz kernel in the distance;
\item \lean{smoothness_order_is_infty}: upstream's smoothness order is $C^\infty$, not analytic;
\item \lean{scaled_density}, \lean{attained}: the hypothesis \lean{DensityRho ρ} is satisfiable, and the bound of D2 is
  attained, so D2 is an optimality statement, not just an inequality (\S\ref{ch2:sec:d2});
\item \lean{unscaled_false}: the naive unscaled statement is false (\S\ref{ch2:sec:d2});
\item \lean{latticeEnergy_riesz_lt_top}: D1 is not $\infty\le\infty$ for $s>2$ (\S\ref{ch2:sec:d1b});
\item \lean{rho_one}, \lean{densityRho_one_iff}: D2 specialises correctly to $\rho=1$;
\item \lean{vacancy_not_isometric_image}: the second minimizer is genuinely different (\S\ref{ch2:sec:d3}).
\end{itemize}

% =====================================================================================================
\section{Exercises}\label{ch2:sec:ex}

Exercises marked \textbf{(Lean)} give a Lean statement with a \lean{sorry} placeholder for you to replace by a proof.
These statements were written for this book and have \emph{not} been compiled by us; part of the exercise is to make
them typecheck (they assume \lean{open OAI.AtomicTriangular Filter Topology} and the namespaces of the files above). A
solution counts only if it passes the axiom gate of \S\ref{ch2:sec:audit}: compile, then \lean{#print axioms}, with
footprint inside the whitelist and no \lean{sorryAx}.

\begin{exercise}[Pen and paper]\label{ch2:ex:neigh}
Using Lemma~\ref{ch2:lem:norm}, find the smallest nonzero value of $\|a\|^2$ for $a\in A$ and all lattice points that
attain it. Show that there are exactly six, and that the minimal distance is $(2/\sqrt3)^{1/2}$.
\end{exercise}

\begin{exercise}[Pen and paper]\label{ch2:ex:diverge}
Show that $\sum_{a\in A\setminus\{0\}}\|a\|^{-s}=\infty$ for $0<s\le2$. \emph{Hint:} bound the number of lattice points
in the annulus $k\le\|a\|<k+1$ from below by a constant times $k$, for large $k$, using the density of $A$.
\end{exercise}

\begin{exercise}[Pen and paper]\label{ch2:ex:poch}
Prove $(-x)^{\downarrow}_r=(-1)^r(x)^{\uparrow}_r$ by induction on $r$, and deduce that $(-1)^r\frac{d^r}{dt^r}t^{-\sigma}>0$
for $\sigma>0$, $t>0$. What goes wrong for $\sigma<0$? Compare with \lean{riesz_neg_two_not_admissible}.
\end{exercise}

\begin{exercise}[Lean]\label{ch2:ex:exp}
Prove that $t\mapsto e^{-t}$ is admissible.
\begin{lstlisting}[language=Lean]
theorem exp_neg_admissible : AdmissiblePotential (fun t => Real.exp (-t)) := by
  sorry
\end{lstlisting}
\emph{Hint:} the $r$-th derivative is $(-1)^re^{-t}$; Mathlib's \lean{iteratedDeriv} lemmas for compositions with
$t\mapsto -t$ may help, or prove $\lean{deriv}^{[r]}$ by induction as in \lean{iterate_deriv_comp_mul}.
\end{exercise}

\begin{exercise}[Pen and paper, then Lean]\label{ch2:ex:const}
Show that the constant $g\equiv1$ is admissible and that $\mathcal{E}_1(A)=\infty$. Then show directly from the
definition that $E_1(C)=\infty$ for every locally finite $C$ of density one ($E_1(C;R)=N_R(C)-1$). Explain why this is
consistent with Theorem~\ref{ch2:thm:upstream}, and why the measure in the Gaussian-mixture representation may charge
$\alpha=0$.
\begin{lstlisting}[language=Lean]
theorem const_admissible : AdmissiblePotential (fun _ => (1 : ℝ)) := by
  sorry
\end{lstlisting}
\end{exercise}

\begin{exercise}[Pen and paper]\label{ch2:ex:homog}
Prove the homogeneity identity $\mathcal{E}_{g_s(\rho^{-1}\cdot)}(A)=\rho^{s/2}\mathcal{E}_{g_s}(A)$ and deduce, for
$s>2$, that the minimal Riesz energy at density $\rho$ is a strictly increasing function of $\rho$. Why does
Proposition~\ref{ch2:prop:unscaled} follow?
\end{exercise}

\begin{exercise}[Pen and paper]\label{ch2:ex:finite}
Let $F\subseteq A$ be finite. Adapt the proof of Theorem~\ref{ch2:thm:vac} to show that $A\setminus F$ has density one and
$E_g(A\setminus F)=\mathcal{E}_g(A)$ for every admissible $g$. Is $A\setminus F$ ever an isometric image of $A$ when
$F\neq\emptyset$?
\end{exercise}

\begin{exercise}[Pen and paper]\label{ch2:ex:add}
Let $s>2$ and $x\notin A$. Show that $A\cup\{x\}$ has density one and that $E_{g_s}(A\cup\{x\})=\mathcal{E}_{g_s}(A)$.
\emph{Hint:} the extra pairs contribute $2\sum_{a\in A\cap B_R}\|x-a\|^{-s}$, which stays bounded as $R\to\infty$; compare
with D1b. For $s\le2$ the conclusion still holds, trivially: both sides are $\infty$. Which step of the finiteness
argument above fails there?
\end{exercise}

\begin{exercise}[Pen and paper]\label{ch2:ex:lf}
Let $C=\{(1-\tfrac1n,0):n\ge1\}\cup\{(2,0)\}$. Show that $C$ is not locally finite, and that there is no $r<1$ with
$C\cap B^\circ_1=C\cap B_r$. Which line of the proof of Lemma~\ref{ch2:lem:sets} fails? What does the Lean convention
``\lean{diskPoints} is empty when the set is infinite'' give for $N_R(C)$, $R\ge1$?
\end{exercise}

\begin{exercise}[Lean]\label{ch2:ex:mono}
Prove that disk counts are monotone in the radius for locally finite configurations.
\begin{lstlisting}[language=Lean]
theorem diskCount_mono {C : Set Plane} (hC : LocallyFinite C) {R R' : ℝ} (h : R ≤ R') :
    diskCount C R ≤ diskCount C R' := by
  sorry
\end{lstlisting}
Then show by an example that the hypothesis \lean{hC} cannot be removed. (Proved in \lean{OpenDisk.lean} as
\lean{diskCount_mono}; use it as the solution key.)
\end{exercise}

\begin{exercise}[Lean]\label{ch2:ex:repar}
Formalise Lemma~\ref{ch2:lem:repar}.
\begin{lstlisting}[language=Lean]
theorem liminf_le_liminf_of_tendsto (u : ℝ → ENNReal) (φ : ℝ → ℝ)
    (hφ : Tendsto φ atTop atTop) :
    liminf u atTop ≤ liminf (u ∘ φ) atTop := by
  sorry
\end{lstlisting}
\emph{Hint:} $\liminf$ along $\lean{atTop}$ of $u\circ\varphi$ is $\liminf u$ along the filter \lean{map φ atTop}, and
\lean{map φ atTop ≤ atTop}. (Proved in \lean{OpenDisk.lean} as \lean{liminf_le_liminf_comp}; use it as the solution
key.)
\end{exercise}

\begin{exercise}[Lean]\label{ch2:ex:vacrho}
Generalise the vacancy theorem to every density: show that $(\rho^{-1/2}A)\setminus\{0\}$ is universally optimal at
density $\rho$.
\begin{lstlisting}[language=Lean]
theorem vacancy_rho_universally_optimal {ρ : ℝ} (hρ : 0 < ρ) :
    TriangularUniversalOptimality.UniversallyOptimal ρ (((Real.sqrt ρ)⁻¹ • A) \ {0}) := by
  sorry
\end{lstlisting}
\emph{Hint:} either redo the proof of \lean{vacancy_energy} at density $\rho$, or show
$((\rho^{-1/2})A)\setminus\{0\}=\rho^{-1/2}(A\setminus\{0\})$ and use the scaling lemmas of \lean{AnyDensity.lean}.
\end{exercise}

\begin{exercise}[Pen and paper]\label{ch2:ex:liminfmul}
Let $u,v:\R\to[0,\infty]$ with $u(R)\to1$. Prove $\liminf uv=\liminf v$ (both inequalities), handling the case
$\liminf v=\infty$. Where does the proof of Theorem~\ref{ch2:thm:vac} need only one of the two inequalities?
\end{exercise}

\begin{exercise}[Discussion]\label{ch2:ex:lim}
Definition~\ref{ch2:def:uo} uses the lower energy. Write down the stronger CKMRV-style definition in which the limit
$\lim_RE_g(C;R)$ must exist. Does the vacancy configuration satisfy it? (You may use, without proof, that
$\lim_RE_g(A;R)$ exists for the potentials in question; state clearly that this is an assumption.) Which of the Lean
theorems of this chapter would need new proofs under the stronger definition?
\end{exercise}

\begin{exercise}[Pen and paper]\label{ch2:ex:d4count}
Show that for the triangular lattice the sequence $N_R(A)-N^\circ_R(A)$ is unbounded as $R\to\infty$ (the boundary
circle can carry many lattice points), and nevertheless both $N_R(A)/(\pi R^2)$ and $N^\circ_R(A)/(\pi R^2)$ tend to $1$.
Why does this not contradict Theorem~\ref{ch2:thm:d4}? \emph{Hint:} the number of representations of an integer by
$j^2+jk+k^2$ is unbounded.
\end{exercise}
